% =====================================================================
% Section 09, Phase 2 introduction
% =====================================================================
\makesectioncover{Phase 2 \textbullet\ The State's Answer}

\section{Phase 2 \textbullet\ What the State Built Since 2014}
\label{sec:phase2-intro}

\subsection{The question Phase 2 is asking}

Phase 1 ended with four numbered gaps. Phase 2 asks whether the
\stat{\$10 billion} Cairo has spent on new transport infrastructure
since 2014 closed any of them.

The state-built infrastructure of the last decade is, in three pieces:

\begin{itemize}
 \item \kw{Cairo Metro Line 3}, Ataba to Adly Mansour in the east,
 then westwards to Kit Kat. The largest single piece of urban-rail
 investment in Egyptian history.
 \item \kw{Cairo LRT}, a light-rail line from Adly Mansour to the
 New Administrative Capital. Designed to anchor the eastward
 expansion of the city.
 \item \kw{Ring-Road BRT}, a bus-rapid-transit corridor along the
 Cairo Ring Road. The most recent and least-publicised of the
 three.
\end{itemize}

Phase 2 collects each of those three networks as data, integrates
them with Phase 1's measured reality, and tests three formal hypotheses
about whether the new infrastructure aligns with where Cairenes
actually live and travel.

\subsection{The Phase 2 pipeline}

Phase 2 also splits across two notebooks:

\begin{itemize}
 \item \fn{Phase2/Phase2\_Scraping\_Cleaning.ipynb},
 scrapes \kw{eight sources} (S1, S3--S8) from Wikipedia, Google
 Maps, OpenStreetMap, citypopulation.de, the GTFS feed, and the
 World Bank API; cleans each source independently; integrates
 them through a four-stage matching pipeline that uses
 \kw{KNN spatial}, \kw{OSM cross-verification}, \kw{RapidFuzz
 fuzzy text} and \kw{SBERT multilingual semantic} matching.
 \item \fn{Phase2/Phase2\_Analysis\_Hypothesis.ipynb}, answers
 \kw{thirteen research questions} (Q13--Q25) and tests
 \kw{three formal hypotheses} (H1, H2, H3). Produces 117+
 Plotly visualisations, three hero maps, and a final
 K-Means market segmentation.
\end{itemize}

\subsection{The two hero maps}

Two visualisations are referenced repeatedly across Phase 2 and are
worth introducing here. Both live as standalone interactive HTML maps
in \fn{Phase2/}; in this report they appear as static PNG renders.

\figitem{hero/two_cairos_map.png}{1.0}{Hero map ``The Two Cairos''. The
formal new modes (Metro L3, LRT, BRT) are layered alongside the Phase
1 informal terminals. The visual signature of the project: two
overlapping but operationally disconnected systems.}{fig:hero-twocairos}

\figitem{hero/headline_coverage_need_map.png}{1.0}{Hero map: the
coverage-need synthesis. Districts coloured by the residual
under-coverage from Q22, with new-mode infrastructure layered on
top. The coral cluster, Imbaba, Shubra, Matariyya, is the
demand the new infrastructure does not visit.}{fig:hero-coverageneed}

These two figures are the visual TL;DR of the entire project. Every
question and hypothesis in Phase 2 either feeds into them or
quantifies what they show.
